\section{Розділ~\ref{sec:dishbrain}. DishBrain}

Будову стенда й результати гри взято з однієї експериментальної роботи та її продовження; формули шуму і дрейфу до добрих налаштувань виведено в розділі й перевірено моделлю книги, а механізм навчання в чашці не встановлено.

{\footnotesize
\setlength{\tabcolsep}{3pt}\renewcommand{\arraystretch}{1.15}
\begin{longtable}{>{\raggedright}p{0.5cm}>{\raggedright}p{4.0cm}>{\raggedright}p{5.6cm}>{\raggedright}p{2.3cm}>{\raggedright\arraybackslash}p{1.7cm}}
\toprule
№ & Твердження & Дослід і що виміряно & Джерело & Статус \\
\midrule\endfirsthead
\toprule
№ & Твердження & Дослід і що виміряно & Джерело & Статус \\
\midrule\endhead
1 & Стенд: близько $800\,000$ клітин кори миші або близько $10^6$ клітин людини на чипі; $26\,000$ електродів на $8$~мм$^2$, запис до $1024$, стимуляція через $8$ & Методи роботи: чип MaxOne, $26\,000$ платинових електродів на $8$~мм$^2$, запис $1024$ каналів по $20\,000$ відліків на секунду; стимулювати технічно можна до $32$ електродів, незалежно вдалося $8$; мишачі культури використовували приблизно з $10$~діб & \cite{Kagan2022} & виміряно \\
2 & Петля з тактом $10$~мс: імпульси рухових ділянок рухають ракетку (рис.~\ref{fig:dishloop}) & Імпульси підраховують вікнами по $10$~мс; затримка від імпульсу до відповідного стимулу близько $5$~мс, її задає буфер апаратури & \cite{Kagan2022} & виміряно \\
3 & Вхід: код місця (який із $8$ електродів) і частота $4$--$40$~імп/с & В основному досліді частота стимуляції зростає лінійно від $4$~Гц, коли м'яч біля дальньої стінки, до $40$~Гц біля ракетки; амплітуда імпульсів м'яча --- $75$~мВ; імпульси двофазні прямокутні & \cite{Kagan2022} & виміряно \\
4 & Вихід $\Delta p=c\,(n_\uparrow-n_\downarrow)$~\eqref{eq:dishreadout}, шум лічби за вікно $1/\sqrt{RT}$ & Правило різниці двох ділянок описано в роботі (з вагами ділянок за активністю), точної формули не дано. Оцінка шуму --- наслідок пуассонівської лічби, виведена в розділі & \cite{Kagan2022}; висновок у розділі & теорія \\
5 & Учитель: після удару --- передбачуваний стимул $75$~мВ, $100$~імп/с, $0.1$~с; після промаху --- непередбачуваний, $150$~мВ, $5$~імп/с, $4$~с у випадкові місця й моменти, потім $4$~с паузи & Методи: після удару $75$~мВ, $100$~Гц, $100$~мс на всі $8$ електродів одразу; після промаху $150$~мВ, $5$~Гц у випадкові електроди у випадкові моменти впродовж $4$~с, потім $4$~с без стимуляції. Дофаміну в культурі немає & \cite{Kagan2022} & виміряно \\
6 & Мережа діє так, щоб її вхід став передбачуваним & Принцип вільної енергії --- теоретична пропозиція, з якої автори виводили протокол; дослід із нею узгоджується, але не відрізняє її від простіших механізмів (рядки~7--8) & \cite{Friston2010,Kagan2022} & гіпотеза \\
7 & Якщо невдача струшує мережу, а вдача ні, час у налаштуванні $\rho\propto1/P_{\text{пр}}$~\eqref{eq:dishdensity} (твердження~\ref{prop:dishscramble}) & Випливає з балансу потоків у марковському ланцюгу із симетричними поштовхами, виведено в розділі. Прямої перевірки на культурі немає & висновок у розділі & теорія \\
8 & Модель «перемішування при промаху» навчається: частка ударів $0.21\to0.38$; за поштовху після кожного розіграшу $\approx0.12$, без поштовхів $0.19$ (рис.~\ref{fig:dishmodel}) & Розрахунок, скрипт \texttt{models/dishbrain\_model.py}, $40$ моделей-культур по $2000$ розіграшів: $0.21\to0.38$; $0.13\to0.11$; $0.19\to0.19$ & --- & модель \\
9 & Серії відбитих м'ячів подовжуються тільки в живих культур у повній петлі; контролі не навчаються & $399$ сеансів по $20$~хв: $9$ культур миші ($101$ сеанс), $11$ людини ($138$), $6$ чипів із середовищем ($80$), $20$ культур у спокої ($42$), $3$ симуляції ($38$). Середня довжина розіграшу за останні $15$~хв більша, ніж за перші $5$, тільки в живих культур; у клітин, що не є нейронами (HEK293T), і в чипів із середовищем ефекту немає & \cite{Kagan2022} & виміряно \\
10 & Значуще зростання за сеанс --- тільки за повного зворотного зв'язку; коли замість стимулу тиша, наприкінці сеансу гра краща, ніж без зворотного зв'язку, але ефект менший & $486$ сеансів за $3$ дні: «стимул» ($n=27$), «тиша» ($17$), «без зворотного зв'язку» ($15$). Зростання в часі значуще тільки в умові «стимул»; наприкінці сеансу умова «тиша» перевершувала «без зворотного зв'язку» з меншим ефектом & \cite{Kagan2022} & виміряно \\
11 & Ефект невеликий; чи навчається мережа надовго --- відкрите питання & Усередині сеансу покращення надійне, але, за словами самих авторів, стійкого навчання між сеансами впродовж кількох днів не спостерігали & \cite{Kagan2022} & виміряно \\
12 & Під час гри мережа наближається до критичного режиму, і що ближче, то краща гра & Аналіз лавин активності тих самих культур: структурований вхід наближає мережу до критичного режиму, якість гри корелює з близькістю до нього; але самої лише критичності без відомостей про наслідки дій для навчання замало & \cite{Habibollahi2023} & виміряно \\
13 & «Розумність» культури не показано & Стаття-відповідь тридцяти дослідників: зміна поведінки в замкненій петлі не доводить ні розуму, ні відчуттів; досліду, який би це показав, немає & \cite{Balci2023} & гіпотеза \\
14 & Запропонований четвертий контроль: передбачуваний стимул після промаху тієї самої тривалості й енергії (підрозділ~\ref{sec:dishspec}) & Такого досліду не поставлено; це пропозиція книги & --- & гіпотеза \\
\bottomrule
\end{longtable}}
